\chapter{Discussion and Analysis}
\label{sec:discussion}

The experimental chapters of this thesis approach political ideology classification from three directions: continued pre-training with metalinguistic auxiliary objectives (Chapter~\ref{sec:continued-pretraining}), supervised fine-tuning of a general-purpose large language model (Chapter~\ref{sec:llm-finetuning}), and causal analysis of the sentiment-to-ideology pathway across human and machine annotators (Chapter~\ref{sec:causal-inference}). Each chapter answers a different question, and their findings speak to one another. The continued pre-training results establish what metalinguistic supervision adds to a contrastive-only objective. The comparative evaluation of specialized transformers against fine-tuned LLMs then places those gains against the headline F1 advantage enjoyed by the larger generative models. The causal analysis, finally, recasts the same classifiers in terms of the textual pathways they rely on, revealing behaviours that the F1-level comparison leaves invisible. This chapter draws those strands together and considers their implications for interpretability, deployment, and the trustworthiness of LLM annotations.

\section{Metalinguistic Feature Integration}
\label{sec:metalinguistic-analysis}

The continued pre-training experiments in Section~\ref{sec:continued-pretraining-results} establish that metalinguistic auxiliary objectives can improve article-level ideology classification, though only under specific configurations. Rather than re-stating the per-configuration F1 numbers, we focus here on what those findings imply about the role of metalinguistic structure in ideological representation.

Three patterns merit particular attention. Theme prediction is the auxiliary objective that delivers the largest and most consistent gains, with the clearest improvements on the historically difficult center and right categories. This is consistent with the cognitive-science framing that topical emphasis is a structural feature of how outlets encode stance, and it suggests that theme complements, rather than substitutes for, the lexical content captured by contrastive pre-training. The benefit of auxiliary objectives is, however, contingent on a sufficiently dense primary contrastive signal: with sparse triplet supervision the auxiliary heads inject more noise than signal. Multi-task pre-training of this form is therefore best understood not as a remedy for weak contrastive supervision but as an amplifier of a sufficiently strong one. Finally, the gains do not stack additively. Combining theme and tone often underperforms theme alone, indicating that the two metalinguistic dimensions overlap in the discriminative information they contribute and that naive summation of auxiliary losses is not the appropriate integration strategy.

Taken together, these results validate the central premise of Chapter~\ref{sec:continued-pretraining}: metalinguistic structure carries ideological signal that is not fully captured by a content-only contrastive objective, and auxiliary metalinguistic supervision is worth including in future continued pre-training work, provided that the primary objective and its supervision density are chosen with care. The gains measured in absolute F1 terms are modest but consistent. Whether they translate into corresponding improvements on the evaluation axes that F1 does not directly measure, namely robustness to shortcut learning and alignment with the human causal profile, is the question the remainder of this chapter takes up.

\input{sections/8_discussion/comparative_analysis.tex}
\section{Causal Analysis: Implications and Interpretation}
\label{sec:causal-discussion}

The causal analysis in Chapter~\ref{sec:causal-inference} reveals a narrow but consistent pattern: the fine-tuned GPT-4o-mini classifier produces the only significant topic-level mediation Total Effect (Donald Trump $\to$ ideology, $-0.249$, with all five paradigms agreeing on the same negative direction) and the only significant community-level ATE in Community~3 (Social Issues, $+0.027$). Its causal profile is the most distinctive of any annotation paradigm, even though every paradigm now agrees with it directionally on the topic-level Trump effect. Human and Fine-tuned RoBERTa annotations produce no significant effects on any Total Effect across the mediation, CATE, or community-level ATE analyses; Fine-tuned RoBERTa's sole significant finding (Hillary Clinton in the multi-treatment analysis) does not replicate under any other paradigm. Two further significant cells appear at the Natural Direct Effect level: baseline GPT-4o-mini in the community $9 \to 11$ pathway and Llama-3.3-70B in the community $11 \to 9$ pathway. As only the ideology labels vary across the datasets, these divergences reflect systematic differences in how each labelling paradigm assigns ideological stance to article text.

\subsection{Shortcut Learning in Fine-Tuned LLMs}

We interpret this divergence as evidence of shortcut learning: fine-tuning on AllSidesV2 amplifies the Trump-sentiment-to-ideology pathway in the model's representation by roughly $3.2\times$ relative to its zero-shot baseline, and produces a single significant community-level coupling on Social Issues sentiment in the rightward direction that does not appear in any other paradigm.

We note that fine-tuned GPT-4o-mini achieves an F1-Macro of 72.48 on the same evaluation set, the highest of any annotator considered here, including human labels. By the usual machine learning yardstick, a classifier with that accuracy is approximating the true label-generating distribution well. That it simultaneously exhibits a causal structure that diverges from the other paradigms in two specific ways (a Total Effect on Donald Trump sentiment that is roughly $3.2\times$ the magnitude of its zero-shot baseline despite full directional agreement across paradigms, and a uniquely significant community-level effect on Social Issues sentiment) is structurally unintuitive and illustrates how causal evaluation complements rather than replaces accuracy metrics, and a single model can pass the F1 test while exhibiting a distinctive causal-structure footprint on the same data.

The pattern of results is consistent with shortcut learning \citep{yuan2024llmsovercomeshortcutlearning, zhou2024explorespuriouscorrelationsconcept}: fine-tuning on AllSidesV2 ideology labels exposes the model to a training distribution in which certain framing patterns are statistically associated with particular ideological labels. The fine-tuned model internalises a stronger coupling between these framings and the target ideology than other paradigms, sharpening shared directional effects into significance and amplifying their magnitudes. As \citet{Ly2024shortcutlearningmedicalAI} show in the medical AI domain, models optimised for in-distribution accuracy can internalise hidden associations that only surface under distributional shift or causal probing. Our mediation analysis serves precisely this diagnostic function.

The results narrow the shortcut-learning evidence to three specific signatures. First, in the topic-level Donald Trump mediation analysis, Fine-tuned GPT-4o-mini's Total Effect ($-0.249$, 95\% CI $[-0.465, -0.033]$) is the only significant cell among the five paradigms, and its magnitude is roughly $3.2\times$ that of the zero-shot GPT-4o-mini baseline ($-0.078$) on the same $N = 233$ articles. Direction agrees with every other paradigm; the distinguishing feature is that fine-tuning bakes a larger effect into the labelling model. The Politics-mediator pathway carries this signal: Fine-tuned GPT-4o-mini's reliable proportion mediated through Politics is $\approx 0.60$ under both percentile contrasts, and the bootstrap NIE is the largest in absolute value across paradigms. Second, in the community-level analysis, Fine-tuned GPT-4o-mini is the only paradigm with a significant ATE on Community~3 (Social Issues, $+0.027$). All five paradigms estimate a positive direction in Community~3, but only Fine-tuned GPT-4o-mini amplifies the shared direction into significance. Third, Fine-tuned GPT-4o-mini is the directional outlier in the community $9 \to 11$ mediation pathway: four of five paradigms estimate a negative Total Effect, and Fine-tuned GPT-4o-mini alone estimates a positive Total Effect ($+0.003$), the cleanest cross-paradigm disagreement in the community-mediation data. Together with the Community~3 amplification, this suggests fine-tuning at the community level (a) sharpens the shared positive Social Issues direction into a significant rightward effect, and (b) flattens the negative Geopolitics direct channel that the other paradigms encode.


\subsection{Implications for LLM Annotation}

A growing line of work explores the potential of LLMs as replacements for human annotators \citep{horton2023large, ma-etal-2025-algorithmic, gao2025LLMscaution}. Our findings caution against this substitution in annotation contexts where downstream causal analyses are planned. A fine-tuned LLM annotator may encode a particular pattern of ideology assignment that does not reflect the full range of human judgments. The Fine-tuned GPT-4o-mini model has internalized a specific annotator profile: one that more strongly associates positive social/cultural sentiment with left-leaning ideology than human annotators do.

If LLM annotations are used as silver-standard training data or as inputs to causal analyses, the spurious pathways learned during fine-tuning will propagate. To make this concrete: consider a researcher analysing whether social-issues framing causally drives article-level ideology classification. If they use Fine-tuned GPT-4o-mini labels as their outcome variable, they detect a significant positive Community~3 effect ($+0.027$, the only paradigm whose interval excludes zero) and conclude that more positive framing of Social Issues sentiment causally shifts ideology classification \emph{rightward}. If they use human labels on exactly the same articles, the effect is small, in the same direction ($+0.011$), and non-significant, and they conclude that Social Issues framing does not robustly drive ideology classification. The two analyses, run on identical text and identical methodology, lead to qualitatively different policy conclusions: the first treats Social Issues framing as a load-bearing causal driver of ideology classification, the second treats it as one of many small directional tendencies that does not clear the significance bar. The fine-tuned LLM does not represent the full complexity of human annotator behaviour; it represents a compressed mapping that amplifies a subset of the directional tendencies present in human labels, sharpening them into significance even when humans do not.

\subsection{Differences Across LLM Families}

The annotators exhibit different causal profiles despite all predicting ideology from the same raw article text, suggesting that architecture and training data differences shape which causal pathways models rely on during inference. Fine-tuned RoBERTa's estimates remain uniformly small and directionally aligned with Human in Community~9 ($-0.015$ vs Human $-0.014$), the cleanest case where the fine-tuned encoder model tracks the human-labelled direct effect. Across the rest of the community-level analysis the new evidence does not clearly distinguish RoBERTa from the LLM-based paradigms in directional terms, since most communities now exhibit full-paradigm directional agreement. We retain the qualitative observation that RoBERTa, despite being fine-tuned on the same AllSidesV2 ideology labels, does not amplify any community-level sentiment coupling into significance the way Fine-tuned GPT-4o-mini does. Two non-exclusive mechanisms could explain this. First, the encoder-only architecture may be structurally less prone to encoding shortcut associations than the generative GPT family. Second, the metalinguistic continued-pretraining objective (Chapter~\ref{sec:continued-pretraining}) may regularise the representation space in a way that resists the kind of sentiment-as-proxy amplification that emerges from naive supervised fine-tuning. We cannot disentangle the two without an ablation we did not run, but the practical implication stands: encoder-based fine-tuning with auxiliary objectives produces uniformly small estimates that do not exhibit the magnitude-amplification footprint of direct fine-tuning of a generative LLM.

Baseline GPT-4o-mini and Fine-tuned GPT-4o-mini share the same underlying architecture, so the divergence between them is attributable entirely to fine-tuning. Baseline GPT-4o-mini produces a significant ATE for Community~9 (Geopolitics, $-0.015$) and a significant Natural Direct Effect for the Community $9 \to 11$ mediation pathway ($-0.093$); its causal profile is therefore not the centrist-bias-only narrative of earlier drafts. Fine-tuned GPT-4o-mini in turn produces a significant ATE for Community~3 (Social Issues, $+0.027$) and the only significant topic-level Trump mediation Total Effect ($-0.249$), with a magnitude roughly $3.2\times$ that of the baseline. The fine-tuning process therefore neither corrects nor cleanly preserves the baseline's causal profile: it amplifies the topic-level Trump-sentiment pathway, sharpens the Social Issues coupling into significance in the rightward direction, and removes the baseline's significant Geopolitics effect.

Llama-3.3-70B shows no significant Total Effect at the community-ATE level despite its pronounced leftward prediction distribution, but its causal profile is not uniformly attenuated. In Community~11 it carries the largest negative magnitude across the five paradigms ($-0.020$), and in Community~9 its point estimate ($-0.014$) sits within rounding distance of the only significant cell (baseline GPT, $-0.015$). In the community $11 \to 9$ mediation analysis, Llama-3.3-70B produces both the largest-magnitude Total Effect ($-0.081$) and the only significant Natural Direct Effect ($-0.057$, 95\% CI $[-0.126, -0.003]$).

The systematic divergences across LLM families, even when controlling for the same input text and outcome task, suggest that ideology annotation by LLMs is not a uniform operation but varies with model-specific feature reliance patterns. This warrants caution in treating any single LLM as a representative annotator.

\subsection{Mediation Analysis as Diagnostic Framework}

A central challenge in deploying LLMs as annotators is that in-distribution F1 alone does not expose why a model produces a given label, and the central claim of this chapter should be stated with that scope. F1 is not the only evaluation tool available: out-of-distribution robustness tests, cross-dataset evaluation, and adversarial probing can sometimes catch the same spurious feature reliance that causal probing catches, and we do not claim otherwise. Causal probing complements robustness benchmarks rather than replacing them, and is most useful in exactly the setting we examine here, where in-distribution accuracy is high but the underlying pathway from sentiment to ideology may diverge from the human one. \citet{moraffah2020causal} argue that causal interpretability, understanding decisions in terms of which inputs caused which outputs, is essential for building trustworthy ML systems, precisely because correlation-based evaluation leaves spurious feature reliance invisible. Our findings illustrate this gap concretely: the fine-tuned model's F1 score (72.48) offered no signal that its ideology predictions amplify the Trump-sentiment-to-ideology pathway by roughly $3.2\times$ relative to the zero-shot baseline, nor that they sharpen the shared Social Issues directional tendency into the only significant rightward community-level coupling.

External causal probing is therefore structurally necessary, not merely complementary to accuracy evaluation. \citet{jin2024largelanguagemodelsinfer} show that LLMs perform near-randomly on tasks requiring pure causal inference from correlational evidence, and \citet{joshi-etal-2024-llms} demonstrate that LLMs are prone to post hoc and ordering fallacies when reasoning about causation. \citet{fu2025correlationcausationanalyzingcausal} further show that standard LLMs rely on spurious correlations in their reasoning processes rather than genuine causal structure, a deficiency that RLVR training partially mitigates but does not eliminate. Taken together, these results imply that LLMs cannot be relied upon to self-report or self-correct the causal basis of their annotations; external causal analysis is required.

Our mediation framework provides one such approach. By holding the input text and causal structure fixed and varying only the annotation paradigm, we isolate differences that are attributable to the annotator rather than the data. This design is consistent with the evaluation principles articulated by \citet{cheng2022evaluation}, who emphasize the need for structured causal evaluation pipelines that go beyond conventional accuracy benchmarks. Applied to annotation auditing, mediation analysis can surface which textual features drive predictions through direct versus mediated pathways, offering a principled basis for trustworthiness assessment that F1 and agreement metrics cannot provide.

